\documentclass[11pt]{article}

\usepackage{fontspec}
\usepackage[spanish,es-noshorthands]{babel}
\usepackage[margin=2.5cm]{geometry}
\usepackage{xurl}
\usepackage{hyperref}
\usepackage{array}
\usepackage{graphicx}
\usepackage{listings}
\usepackage{longtable}
\usepackage{xcolor}

\setmainfont{Latin Modern Roman}
\setmonofont{Latin Modern Mono}
\setlength{\emergencystretch}{4em}
% Permite cortar línea después de un guion bajo (nombres largos de pruebas)
\let\guionbajo\_
\renewcommand{\_}{\guionbajo\allowbreak}

\lstset{
  basicstyle=\ttfamily\small,
  columns=fullflexible,
  keepspaces=true,
  breaklines=true,
  frame=single,
  framesep=4pt,
  xleftmargin=6pt,
  xrightmargin=6pt,
  inputencoding=utf8,
  extendedchars=true,
  literate={á}{{\'a}}1 {é}{{\'e}}1 {í}{{\'i}}1 {ó}{{\'o}}1 {ú}{{\'u}}1 {ñ}{{\~n}}1
}
\lstdefinestyle{java}{
  language=Java,
  keywordstyle=\bfseries,
  commentstyle=\itshape\color{gray!70!black},
  numbers=left,
  numberstyle=\scriptsize\color{gray},
  numbersep=6pt,
  xleftmargin=18pt
}

% Captura: usa capturas/#1 si existe; si no, deja el recuadro
\graphicspath{{capturas/}}
\newcommand{\captura}[2]{%
  \begin{center}
  \IfFileExists{capturas/#1}%
    {\includegraphics[width=\linewidth,height=0.42\textheight,keepaspectratio]{#1}}%
    {\fbox{\parbox[c][3cm][c]{0.9\linewidth}{\centering\small Captura pendiente: \texttt{capturas/#1}}}}\\[0.6em]
  {\small #2}
  \end{center}}

% Encabezado de cada evidencia: número, título corto y puntos
\newcommand{\evidencia}[3]{\subsection*{#1. #2 \hfill {\normalsize\mdseries (#3)}}}

% Nombre de una prueba (se escribe con guiones bajos normales): puede cortarse en cada guion
% bajo y antes de cada mayúscula del camelCase, para que los nombres largos no se salgan
\directlua{dofile("prueba.lua")}
\newcommand{\prueba}[1]{\texttt{\directlua{prueba("\luaescapestring{\detokenize{#1}}")}}}

\newcommand{\repo}{https://github.com/Vethariel/Lab6-ingesoft2}

\hypersetup{
  pdftitle={Laboratorio 6. Testing},
  pdfauthor={Daniel Alonso Gracia Pinto},
  colorlinks=true,
  linkcolor=black,
  urlcolor=blue
}

\title{Laboratorio 6: Testing\\[0.35em]
\large Pruebas unitarias e integración, TDD, cobertura y calidad\\[0.25em]
\normalsize Ingeniería de Software II (2016702)}
\author{Daniel Alonso Gracia Pinto}
\date{Noviembre de 2026\\[0.35em]
\normalsize Grupo 00}

\begin{document}
\maketitle
\vspace{-1.6em}

\begin{center}
\begin{tabular}{@{}>{\raggedright\arraybackslash}p{3.4cm}>{\raggedright\arraybackslash}p{11.6cm}@{}}
\hline
Integrante & Daniel Alonso Gracia Pinto \\
Grupo & 00 \\
Proyecto & \texttt{Prestamos-00} \\
Repositorio & \url{\repo} \\
\hline
\end{tabular}
\end{center}

El repositorio contiene el proyecto completo en \texttt{Prestamos-00/}, junto con el historial de commits (incluidos los de cada ciclo de TDD) y el \texttt{README.md}, que explica cómo ejecutar las pruebas y la cobertura. Todas las capturas son de pantalla completa, de modo que en cada una se ven el nombre del proyecto, \texttt{Prestamos-00}, y la hora del sistema en la esquina inferior derecha.

\section*{Paso 0. Stack y proyecto}

Se eligió Java con Maven, el mismo lenguaje de los laboratorios 4 y 5, porque todas las piezas que pide el laboratorio tienen allí una herramienta madura. En concreto, el stack queda así:

\begin{center}
\begin{tabular}{@{}>{\raggedright\arraybackslash}p{4.6cm}>{\raggedright\arraybackslash}p{6.4cm}>{\raggedright\arraybackslash}p{3.6cm}@{}}
\hline
Pieza & Herramienta & Versión \\
\hline
(1) Framework de pruebas & JUnit 5 (Jupiter) & 5.14.2 \\
(2) Dobles de prueba & Mockito & 5.20.0 \\
(3) Cliente HTTP y JSON & \texttt{java.net.http.HttpClient} (del JDK) y Gson & Gson 2.13.2 \\
(4) Servidor HTTP de prueba & MockWebServer (OkHttp) & 5.3.0 \\
(5) Cobertura & JaCoCo & 0.8.14 \\
\hline
\end{tabular}
\end{center}

El código de producción se divide en dos paquetes. Por un lado, \texttt{co.edu.unal.\allowbreak prestamos.domain} contiene las reglas del negocio y no depende de red, base de datos ni framework. Por otro, \texttt{co.edu.unal.\allowbreak prestamos.data} contiene los adaptadores que hablan con servicios externos. Las pruebas repiten esa misma estructura en \texttt{src/test/java}, así que cada prueba queda junto al paquete que verifica.

\evidencia{E1}{Dependencias y prueba trivial en verde}{0,15}

\captura{e1-a-dependencias.png}{E1 (a). \texttt{pom.xml}: las dependencias, cada una marcada con la pieza del stack que cubre.}

\captura{e1-b-plugins.png}{E1 (b). \texttt{pom.xml}: Surefire, que ejecuta las pruebas, y JaCoCo, que mide la cobertura.}

\captura{e1-c-prueba-trivial.png}{E1 (c). \texttt{mvn test} con la prueba trivial \texttt{EntornoTest} ($1 + 1 = 2$): \texttt{Tests run: 1, Failures: 0} y \texttt{BUILD SUCCESS}. Una vez confirmado que el entorno funciona, la prueba se borró.}

\evidencia{E2}{Estructura de carpetas y README}{0,15}

\captura{e2-estructura-readme.png}{E2. A la izquierda, \texttt{domain} y \texttt{data} tanto en \texttt{src/main} como en \texttt{src/test}; a la derecha, el \texttt{README.md} con la tabla del stack.}

\section*{Parte 1. Pruebas unitarias del dominio}

\texttt{Prestamo} es el agregado del dominio: guarda su propio estado y solo lo cambia cuando se cumplen las reglas. Cuando una regla no se cumple, en cambio, lanza \texttt{IllegalStateException} con el mensaje del enunciado. Además, los 7 días del préstamo, los 7 de la renovación y el máximo de una renovación están en constantes (\texttt{DIAS\_DE\_PRESTAMO}, \texttt{DIAS\_DE\_RENOVACION} y \texttt{MAXIMO\_DE\_RENOVACIONES}), de modo que cada regla tiene un nombre. Por último, las pruebas usan una fecha fija, \texttt{HOY = 2026-11-18}, en lugar de la fecha del sistema; así dan el mismo resultado sin importar el día en que se ejecuten.

\evidencia{E3}{Código de las siete pruebas}{0,4}

Todas las pruebas siguen la estructura AAA: sus bloques están marcados con \texttt{// Arrange}, \texttt{// Act} y \texttt{// Assert}, y su nombre sigue el patrón \texttt{metodo\_escenario\_resultadoEsperado}. Como todas parten de un préstamo pendiente, ese Arrange común está en un método \texttt{@BeforeEach}, que JUnit ejecuta antes de cada prueba.

\begin{center}
\begin{tabular}{@{}>{\raggedright\arraybackslash}p{0.6cm}>{\raggedright\arraybackslash}p{9.2cm}>{\raggedright\arraybackslash}p{4.8cm}@{}}
\hline
 & Prueba & Escenario \\
\hline
-- & \prueba{aprobar_prestamoPendiente_quedaAprobadoConLimiteEnSieteDias} & Modelo del enunciado \\
-- & \prueba{renovar_porSegundaVez_lanzaError} & Modelo del enunciado \\
(a) & \prueba{renovar_prestamoAprobado_extiendeLimiteACatorceDias} & Límite en 2026-12-02 \\
(b) & \prueba{aprobar_prestamoYaAprobado_lanzaError} & ``Solo se aprueba un préstamo pendiente'' \\
(c) & \prueba{renovar_prestamoPendiente_lanzaErrorDePrestamoNoAprobado} & ``Solo se renueva un préstamo aprobado'' \\
(d) & \prueba{devolver_prestamoAprobado_quedaDevuelto} & Estado \texttt{DEVUELTO} \\
(e) & \prueba{devolver_prestamoYaDevuelto_lanzaError} & ``Solo se devuelve un préstamo aprobado'' \\
\hline
\end{tabular}
\end{center}

\captura{e3-a-pruebas.png}{E3 (a). \texttt{PrestamoTest}, líneas 1 a 45: la fecha fija, el \texttt{@BeforeEach} y las dos pruebas modelo.}

\captura{e3-b-pruebas.png}{E3 (b). Líneas 45 a 87: los escenarios (a), (b) y (c), y el inicio de (d).}

\captura{e3-c-pruebas.png}{E3 (c). Líneas 75 a 105: los escenarios (d) y (e).}

\evidencia{E4}{Las siete pruebas en verde}{0,3}

\captura{e4-verde.png}{E4. \texttt{mvn test}: \texttt{PrestamoTest} con \texttt{Tests run: 7, Failures: 0} y \texttt{BUILD SUCCESS}.}

\evidencia{E5}{La prueba que falla con 8 días}{0,3}

\captura{e5-falla-ocho-dias.png}{E5. Con \texttt{DIAS\_DE\_PRESTAMO = 8} fallan dos pruebas: \texttt{expected: <2026-11-25> but was: <2026-11-26>} y \texttt{expected: <2026-12-02> but was: <2026-12-03>}.}

\textbf{Qué detectó la prueba:} \prueba{aprobar_prestamoPendiente_quedaAprobadoConLimiteEnSieteDias} detectó que, al aprobar, la fecha límite quedaba un día más tarde de lo que dice la regla: esperaba el 25 de noviembre (7 días desde el 18), pero obtuvo el 26. De paso, también falló la prueba de renovación, ya que su fecha esperada (14 días) se calcula a partir de la misma aprobación. Es decir, un solo error en la regla rompió todo lo que depende de ella, y las pruebas lo señalaron con el valor exacto. Después de la captura se restauró el valor 7.

\section*{Parte 2. Desarrollo guiado por pruebas (TDD)}

\texttt{CalculadoraMulta} se construyó en cuatro ciclos, uno por regla. En cada ciclo, primero se escribió la prueba y se comprobó que fallaba por la razón esperada (rojo); luego se escribió el mínimo código para que pasara (verde), y por último se mejoró el diseño sin cambiar el comportamiento (refactor). Para contar los días de retraso se usa \texttt{ChronoUnit.DAYS.between}, como pide el enunciado, y las pruebas parten de una fecha fija, \texttt{LIMITE = 2026-11-25}.

\begin{center}
\begin{tabular}{@{}>{\raggedright\arraybackslash}p{1.2cm}>{\raggedright\arraybackslash}p{5.6cm}>{\raggedright\arraybackslash}p{7.8cm}@{}}
\hline
Ciclo & Prueba que lo inicia (rojo) & Código del verde \\
\hline
1 & \prueba{calcular_devueltoAntesDelLimite_noHayMulta} & \texttt{return 0;} \\
2 & \prueba{calcular_unDiaDeRetraso_cobraDosMil} & Contar los días de retraso y multiplicar por 2000 \\
3 & \prueba{calcular_cuarentaDiasDeRetraso_noPasaDelTope} & \texttt{Math.min(..., 50000)} \\
4 & \prueba{calcular_monitorConTresDiasDeRetraso_pagaLaMitad} & Si es monitor, dividir entre 2 después del tope \\
\hline
\end{tabular}
\end{center}

Además de la prueba que inicia cada ciclo, se agregaron los casos límite del enunciado: devuelto exactamente el día límite (\$0), 25 días de retraso (\$50.000, justo el tope) y un monitor con 40 días (\$25.000). Este último es importante porque comprueba que la mitad se aplica \emph{después} del tope; si se aplicara antes, la multa sería $40 \times 1.000 = 40.000$, y la prueba fallaría.

\evidencia{E6}{Rojo y verde del ciclo 1}{0,3}

\captura{e6-a-ciclo1-rojo.png}{E6 (a). Rojo: la prueba ya existe, pero \texttt{CalculadoraMulta} todavía no. Como Java es un lenguaje compilado, el rojo aparece como error de compilación: \texttt{cannot find symbol: class CalculadoraMulta}.}

\captura{e6-b-ciclo1-verde.png}{E6 (b). Verde: \texttt{CalculadoraMulta} con \texttt{return 0;}, que es lo mínimo para pasar. \texttt{Tests run: 8, Failures: 0}.}

\evidencia{E7}{Rojo y verde de los ciclos 2, 3 y 4}{0,45}

\captura{e7-a-ciclo2-rojo.png}{E7 (a). Ciclo 2, rojo: \prueba{calcular_unDiaDeRetraso_cobraDosMil} falla con \texttt{expected: <2000> but was: <0>}, porque el código todavía devuelve 0.}

\captura{e7-b-ciclo2-verde.png}{E7 (b). Ciclo 2, verde: la multa pasa a ser los días de retraso por 2000. \texttt{Tests run: 10, Failures: 0}.}

\captura{e7-c-ciclo3-rojo.png}{E7 (c). Ciclo 3, rojo: \prueba{calcular_cuarentaDiasDeRetraso_noPasaDelTope} falla con \texttt{expected: <50000> but was: <80000>}, ya que aún no hay tope. En cambio, la de 25 días ya pasa.}

\captura{e7-d-ciclo3-verde.png}{E7 (d). Ciclo 3, verde: \texttt{Math.min(diasDeRetraso * MULTA\_POR\_DIA, 50000)}. \texttt{Tests run: 12, Failures: 0}.}

\captura{e7-e-ciclo4-rojo.png}{E7 (e). Ciclo 4, rojo: las dos pruebas de monitor fallan con \texttt{expected: <3000> but was: <6000>} y \texttt{expected: <25000> but was: <50000>}, porque el código aún no distingue a los monitores.}

\captura{e7-f-ciclo4-verde.png}{E7 (f). Ciclo 4, verde: si es monitor, la multa con tope se divide entre 2. \texttt{Tests run: 14, Failures: 0}.}

\evidencia{E8}{Historial de commits}{0,25}

\captura{e8-git-log.png}{E8. \texttt{git log --oneline}: el ciclo 1 en verde y, después, los ciclos 2, 3 y 4 con sus commits de rojo, verde y refactor, en ese orden.}

El historial deja ver la disciplina de TDD: cada commit de rojo contiene la prueba nueva fallando, mientras que ningún commit de verde cambia las pruebas, solo el código. Los refactors, por su parte, se limitaron a mejorar nombres y estructura. En el ciclo 2 se creó la constante \texttt{MULTA\_POR\_DIA} y las pruebas pasaron a compartir una sola calculadora; en el ciclo 3 se creó \texttt{MULTA\_MAXIMA}, y en el ciclo 4, \texttt{SIN\_MULTA}, \texttt{DIVISOR\_MONITOR} y el método \texttt{multaConTope}.

\evidencia{E9}{Código final de \texttt{CalculadoraMulta}}{0,2}

\captura{e9-calculadora-final.png}{E9. \texttt{CalculadoraMulta} después del último refactor, sin números mágicos.}

\lstinputlisting[style=java,firstline=6]{../Prestamos-00/src/main/java/co/edu/unal/prestamos/domain/CalculadoraMulta.java}

\textbf{Cómo influyeron las pruebas en el diseño.} Como la prueba se escribía primero, \texttt{calcular} tuvo que recibir las dos fechas como parámetros en lugar de leer la fecha del sistema; de lo contrario, la prueba no podría fijar el resultado. En consecuencia, la calculadora quedó como una función pura y sin estado, fácil de probar. Del mismo modo, las pruebas de los casos límite fijaron dónde va el \texttt{<=} y dónde se aplica el tope, y la del monitor con 40 días fijó el orden de las reglas. Finalmente, como todas las pruebas estaban en verde, el refactor pudo dar un nombre a cada regla sin temor a romper nada.

\section*{Parte 3. Dobles de prueba}

En \texttt{domain} se definieron tres contratos como interfaces de Java (\texttt{DisponibilidadEquipos}, \texttt{PrestamoRepository} y \texttt{Notificador}), el error \texttt{EquipoNoDisponible} y el caso de uso \texttt{SolicitarPrestamo}, que recibe los tres contratos por constructor. Gracias a eso, las pruebas pueden entregarle dobles en lugar de implementaciones reales. Concretamente, un \texttt{@BeforeEach} crea los tres mocks con \texttt{mock(...)} de Mockito; luego cada prueba configura con \texttt{when} o \texttt{doThrow} lo que necesita y, al final, comprueba con \texttt{verify} qué llamadas se hicieron.

\evidencia{E10}{Código de las tres pruebas con dobles}{0,3}

\begin{center}
\begin{tabular}{@{}>{\raggedright\arraybackslash}p{5.6cm}>{\raggedright\arraybackslash}p{4.2cm}>{\raggedright\arraybackslash}p{4.8cm}@{}}
\hline
Prueba & Configuración & Verificación \\
\hline
\prueba{ejecutar_equipoDisponible_guardaYNotifica} & \texttt{estaDisponible(12)} $\to$ \texttt{true}; \texttt{siguienteId()} $\to$ \texttt{50} & Estado \texttt{PENDIENTE}; \texttt{guardar(p)} una vez; \texttt{notificar("1023", "Solicitud 50 registrada")} \\
(a) \prueba{ejecutar_equipoNoDisponible_noGuardaNiNotifica} & \texttt{estaDisponible(12)} $\to$ \texttt{false} & Lanza \texttt{EquipoNoDisponible}; \texttt{guardar} nunca (\texttt{never()}); el notificador sin llamadas (\texttt{verifyNoInteractions}) \\
(b) \prueba{ejecutar_falloAlGuardar_noNotifica} & \texttt{guardar(...)} lanza \texttt{RuntimeException("BD caída")} & El error se propaga con su mensaje; el notificador sin llamadas \\
\hline
\end{tabular}
\end{center}

\captura{e10-a-dobles.png}{E10 (a). \texttt{SolicitarPrestamoTest}: los imports de Mockito, la creación de los tres mocks y la primera prueba.}

\captura{e10-b-dobles.png}{E10 (b). Las pruebas (a) y (b), con \texttt{doThrow}, \texttt{never()} y \texttt{verifyNoInteractions}.}

\evidencia{E11}{Las tres pruebas en verde}{0,3}

\captura{e11-verde.png}{E11. \texttt{mvn test}: \texttt{SolicitarPrestamoTest} con \texttt{Tests run: 3, Failures: 0}, junto con las demás pruebas del dominio (17 en total).}

\evidencia{E12}{Qué verifica un mock y por qué se pueden usar dobles}{0,2}

\textbf{¿Qué verifica un mock que no puede verificar una aserción sobre el valor de retorno?} Verifica el \emph{comportamiento} del caso de uso hacia afuera, es decir, a quién llamó, con qué argumentos, cuántas veces y a quién no llamó. Esto importa porque lo esencial de \texttt{SolicitarPrestamo} son sus efectos, que el préstamo se guarde y que el estudiante reciba el aviso, y esos efectos no aparecen en el valor de retorno. Por ejemplo, la primera prueba podría recibir un préstamo correcto aunque nunca se hubiera guardado, y solo \texttt{verify(repositorio).guardar(p)} lo detectaría. En las pruebas (a) y (b), además, ni siquiera hay valor de retorno, porque \texttt{ejecutar} lanza un error; por lo tanto, lo único que se puede comprobar es lo que \emph{no} pasó: que no se guardó (\texttt{never()}) y que no se notificó (\texttt{verifyNoInteractions}). Así se descarta, por ejemplo, avisarle al estudiante de una solicitud que la base de datos nunca registró.

\textbf{¿Por qué es posible usar dobles aquí?} Porque el diseño aplica el principio de inversión de dependencias (DIP): \texttt{SolicitarPrestamo} no depende de clases concretas, como un cliente HTTP o una base de datos, sino de tres interfaces que el propio dominio define y que recibe por constructor. En términos de arquitectura hexagonal, esas interfaces son los \emph{puertos} del dominio, mientras que sus implementaciones (el adaptador HTTP de la Parte 4, un repositorio con base de datos o un servicio de correo) son los \emph{adaptadores}, que viven afuera, en \texttt{data}. Como el caso de uso no sabe qué adaptador tiene detrás, la prueba puede conectar un mock en cada puerto. En cambio, si \texttt{SolicitarPrestamo} creara por dentro un \texttt{new DisponibilidadEquiposHttp(...)}, no habría forma de reemplazarlo, y probarlo exigiría la red.

\section*{Parte 4. Pruebas de integración con un servidor HTTP de prueba}

\texttt{DisponibilidadEquiposHttp}, en \texttt{data}, es el adaptador que cumple el contrato \texttt{DisponibilidadEquipos}. Para hacerlo, consulta \texttt{GET <urlBase>/api/v1/equipos/<id>} con \texttt{HttpClient}, convierte el JSON de la respuesta con Gson y recibe la URL base por constructor, de modo que la prueba puede apuntarlo a su servidor local. Su comportamiento depende de la respuesta: un 200 con \texttt{"estado": "disponible"} devuelve verdadero; un 200 con otro estado o un 404 devuelven falso; cualquier otro código lanza \texttt{ServicioNoDisponible("HTTP <código>")} y, por último, un error de red lanza \texttt{ServicioNoDisponible} con el mensaje del error.

Cabe aclarar que \texttt{ServicioNoDisponible} quedó en \texttt{domain} y no en \texttt{data}. La razón es que forma parte del contrato del puerto: así, quien use el caso de uso puede atrapar ese error sin depender de la capa de infraestructura.

A diferencia de la Parte 3, estas pruebas no reemplazan el adaptador ni el cliente HTTP, sino que usan \textbf{MockWebServer}, un servidor HTTP real que escucha en un puerto local. Cada prueba lo inicia en \texttt{@BeforeEach}, le indica qué respuesta dar y lo cierra en \texttt{@AfterEach}. De esta forma, la petición sale por el \texttt{HttpClient} real, viaja por la red local y la respuesta pasa por Gson, igual que en producción.

\evidencia{E13}{Pruebas (a), (b) y (c) con MockWebServer}{0,35}

\captura{e13-a-servidor.png}{E13 (a). \texttt{DisponibilidadEquiposHttpTest}: el servidor \texttt{MockWebServer}, el método \texttt{responder}, que encola un \texttt{MockResponse}, y la primera prueba, que con \texttt{takeRequest()} verifica el método \texttt{GET} y la ruta \texttt{/api/v1/equipos/12}.}

\captura{e13-b-casos.png}{E13 (b). Los casos (a) 200 con \texttt{"prestado"} $\to$ falso, (b) 404 $\to$ falso y (c) 500 $\to$ \texttt{ServicioNoDisponible("HTTP 500")}.}

\evidencia{E14}{Fake en memoria y prueba de extremo a extremo}{0,3}

\captura{e14-a-fake.png}{E14 (a). \texttt{PrestamoRepositoryEnMemoria}: el fake es solo una lista y un contador. Sin embargo, a diferencia de un mock, funciona de verdad, pues guarda y numera los préstamos.}

\captura{e14-b-extremo.png}{E14 (b). \texttt{SolicitarPrestamoIntegracionTest}: \texttt{SolicitarPrestamo} con el adaptador HTTP real contra el servidor de prueba, el fake y un notificador mock. Al final, verifica que el préstamo quedó guardado en el fake con id 1.}

\evidencia{E15}{Todas las pruebas de \texttt{data} en verde}{0,2}

\captura{e15-data-verde.png}{E15. \texttt{mvn test -Dtest='co.edu.unal.prestamos.data.**'}: \texttt{Tests run: 5, Failures: 0}.}

\evidencia{E16}{Qué detecta esta prueba que los mocks nunca detectarían}{0,15}

En la Parte 3, \texttt{DisponibilidadEquipos} era un mock que respondía lo que la prueba le indicara; por lo tanto, el adaptador nunca se ejecutaba, y cualquier error en él pasaba inadvertido. La prueba con el servidor, en cambio, recorre el camino completo. Tres ejemplos concretos lo muestran:

\begin{enumerate}
  \item \textbf{La ruta.} Si el adaptador armara mal la URL, el servidor real respondería 404 y, en consecuencia, el adaptador diría ``no disponible'' para todos los equipos. Esto pasaría, por ejemplo, con \texttt{/api/v1/equipo/12} (en singular) o con una barra doble, si \texttt{urlBase} terminara en \texttt{/}. De hecho, este último caso no es hipotético: \texttt{servidor.url("/")} devuelve la URL con barra final, y por eso el constructor la quita. La primera prueba lo cubre, ya que lee la petición con \texttt{takeRequest()} y compara la ruta exacta.
  \item \textbf{El JSON.} Si la API devolviera \texttt{"status"} en lugar de \texttt{"estado"}, o \texttt{"Disponible"} con mayúscula, Gson dejaría el campo en \texttt{null} o con otro valor, y el adaptador devolvería falso sin lanzar ningún error. Como con mocks nunca se convierte un JSON, este error solo aparece cuando una respuesta real pasa por Gson.
  \item \textbf{Los códigos HTTP.} Confundir el 404 (``no existe'', que se responde con falso) con un error del servidor (que lanza \texttt{ServicioNoDisponible}) cambia lo que ve el estudiante. Mientras el mock del puerto solo devuelve \texttt{true} o \texttt{false} y nunca ve un código, las pruebas (b) y (c) sí lo ven. Además, el tercer mutante de la Parte 5 (404 $\to$ 400) lo confirma, porque lo mata precisamente la prueba (b).
\end{enumerate}

\section*{Parte 5. Cobertura y calidad de las pruebas}

\textbf{Medición inicial.} Con las pruebas de las partes 1 a 4, JaCoCo reportó 94,7~\% de líneas en \texttt{domain} (54 de 57) y 76,0~\% en \texttt{data} (19 de 25). Es decir, \texttt{domain} ya cumplía la meta, pero \texttt{data} no. Según el reporte, faltaban por ejecutar los dos \texttt{catch} del adaptador y, en el dominio, el constructor de \texttt{ServicioNoDisponible} que recibe la causa y \texttt{getEstudianteId()}.

\textbf{Pruebas agregadas.} Para cubrir el error de red se agregó \prueba{estaDisponible_servidorApagado_lanzaServicioNoDisponibleConElErrorDeRed}. Esta prueba apaga el servidor antes de la petición, de modo que la conexión es rechazada y el adaptador debe convertir el \texttt{ConnectException} en \texttt{ServicioNoDisponible}. Además, la prueba de extremo a extremo ahora verifica también el \texttt{estudianteId}. Con esto, la cobertura llegó a \textbf{100~\% en \texttt{domain} (57 de 57) y 88~\% en \texttt{data} (22 de 25)}.

Las 3 líneas que siguen sin cubrir son el \texttt{catch (InterruptedException)}, que solo se ejecuta si otro hilo interrumpe al que espera la respuesta. Provocarlo exigiría coordinar dos hilos con tiempos, lo que daría una prueba frágil; por eso no se forzó, más aún cuando la meta ya se cumple y ese bloque solo aplica el patrón estándar de restaurar la marca de interrupción.

Por último, para que la meta no dependa de revisar el reporte a mano, el \texttt{pom.xml} incluye la regla \texttt{check} de JaCoCo. Gracias a ella, \texttt{mvn verify} falla si \texttt{domain} baja de 90~\% o \texttt{data} de 80~\% de líneas.

\evidencia{E17}{Resumen de cobertura que cumple la meta}{0,2}

\captura{e17-cobertura.png}{E17. \texttt{mvn verify}: las 23 pruebas en verde, el reporte y la comprobación \texttt{jacoco:check} con \texttt{All coverage checks have been met.}}

\begin{center}
\begin{tabular}{@{}>{\raggedright\arraybackslash}p{4.4cm}>{\raggedright\arraybackslash}p{3.4cm}>{\raggedright\arraybackslash}p{3.4cm}>{\raggedright\arraybackslash}p{3.2cm}@{}}
\hline
Paquete & Líneas & Ramas & Meta de líneas \\
\hline
\texttt{domain} & 57 de 57: 100~\% & 14 de 14: 100~\% & 90~\%: cumple \\
\texttt{data} & 22 de 25: 88~\% & 6 de 8: 75~\% & 80~\%: cumple \\
\hline
\end{tabular}
\end{center}

\evidencia{E18}{Reporte HTML en el navegador}{0,15}

\captura{e18-a-reporte.png}{E18 (a). \texttt{target/site/jacoco/index.html} en Firefox. Las líneas se leen en las columnas \emph{Missed} y \emph{Lines}: 3 sin ejecutar de 25 en \texttt{data} y 0 de 57 en \texttt{domain}. La columna \emph{Cov.} de la izquierda, en cambio, mide instrucciones y no líneas.}

\captura{e18-b-detalle.png}{E18 (b). El código de \texttt{DisponibilidadEquiposHttp} en el reporte: en verde lo ejecutado, en amarillo las ramas recorridas solo en parte y en rojo el \texttt{catch (InterruptedException)}, que es lo único sin ejecutar.}

\evidencia{E19}{Mutación manual}{0,2}

Cada mutante se introdujo por separado; luego se ejecutaron las 23 pruebas y, al terminar, se restauró el código con \texttt{git checkout}. Hay que aclarar que, en el código, la condición \texttt{renovaciones < 1} está escrita como \texttt{renovaciones >= MAXIMO\_DE\_RENOVACIONES}, con la constante en 1. Por eso, el primer mutante cambia esa constante a 2, lo que equivale a \texttt{renovaciones < 2}.

\begin{center}
\begin{tabular}{@{}>{\raggedright\arraybackslash}p{3.8cm}>{\raggedright\arraybackslash}p{1.8cm}>{\raggedright\arraybackslash}p{4.4cm}>{\raggedright\arraybackslash}p{4.6cm}@{}}
\hline
Mutante & Resultado & Prueba que lo mata & Mensaje \\
\hline
\texttt{renovar}: \texttt{renovaciones < 1} $\to$ \texttt{< 2} & Muerto & \prueba{renovar_porSegundaVez_lanzaError} & \texttt{Expected IllegalStateException to be thrown, but nothing was thrown} \\
\texttt{CalculadoraMulta}: tope \$50.000 $\to$ \$60.000 & Muerto & \prueba{calcular_cuarentaDiasDeRetraso_noPasaDelTope} y \prueba{calcular_monitorConCuarentaDiasDeRetraso_pagaLaMitadDelTope} & \texttt{expected: <50000> but was: <60000>} y \texttt{expected: <25000> but was: <30000>} \\
Adaptador: \texttt{404} $\to$ \texttt{400} & Muerto & \prueba{estaDisponible_equipoInexistente404_devuelveFalso} & \texttt{ServicioNoDisponible: HTTP 404} \\
\hline
\end{tabular}
\end{center}

Como ningún mutante sobrevivió, no hubo que escribir pruebas nuevas para matarlos. Aun así, el segundo muestra lo cerca que estuvo de sobrevivir. La prueba de 25 días también ejecuta la línea del tope, pero no lo detecta, porque $25 \times 2.000 = 50.000$ no supera ninguno de los dos topes. En realidad, lo matan solo las pruebas de 40 días, que se habían escrito en la Parte 2 como casos límite.

\captura{e19-mutante-tope.png}{E19. El mutante del tope: \texttt{MULTA\_MAXIMA = 60\_000} en el editor (marcado como modificado) y las dos pruebas de 40 días fallando. El \texttt{BUILD FAILURE} indica que el mutante murió.}

\evidencia{E20}{Por qué 100~\% de cobertura no garantiza que el código sea correcto}{0,15}

La cobertura mide qué líneas se \emph{ejecutaron} durante las pruebas, pero no qué se \emph{verificó}. En otras palabras, una línea cuenta como cubierta aunque ninguna aserción compruebe su resultado.

El segundo mutante lo demuestra en este mismo proyecto. Si \texttt{CalculadoraMultaTest} tuviera solo las pruebas de 0, 1 y 25 días, se ejecutarían todas las líneas de \texttt{CalculadoraMulta}, incluida la del \texttt{Math.min} con el tope, y el reporte marcaría 100~\%. Sin embargo, con el tope en \$60.000 esas pruebas seguirían en verde, porque ninguna llega a una multa que lo supere. Así, el código estaría mal con 100~\% de cobertura. Lo que mató al mutante, entonces, no fue ejecutar la línea, sino tener una prueba de 40 días con una aserción capaz de distinguir \$50.000 de \$60.000. Llevado al extremo, una prueba sin aserciones también ``cubre'' líneas sin verificar nada.

Además, hay errores que la cobertura no puede ver. Si al código le falta una regla (por ejemplo, el descuento para monitores antes del ciclo 4), no existen líneas que dejar sin cubrir; y si una prueba tiene mal el valor esperado, la cobertura sale igual. Por eso las dos herramientas se complementan: la cobertura sirve para encontrar código que ninguna prueba toca, como el \texttt{catch} de la E18, mientras que la mutación mide si las pruebas que sí lo tocan detectan los errores. En este proyecto, el 100~\% de líneas en \texttt{domain} y los tres mutantes muertos dicen más juntos que cualquiera de los dos por separado.

\section*{Entrega}

El repositorio \url{\repo} contiene:

\begin{description}
  \raggedright
  \item[\texttt{Prestamos-00/}.] El proyecto Maven, con \texttt{domain} y \texttt{data}, sus 23 pruebas y el \texttt{pom.xml}, donde están el stack y las metas de cobertura.
  \item[\texttt{Prestamos-00/README.md}.] La tabla del stack y los comandos: \texttt{mvn test} ejecuta las pruebas y \texttt{mvn verify} mide la cobertura, cuyo reporte queda en \texttt{target/site/jacoco/index.html}. Basta con tener Java 17 o superior y Maven, ya que las dependencias se descargan solas.
  \item[El historial de commits.] Un commit por parte y, en la Parte 2, los commits \texttt{TDD ciclo N: rojo}, \texttt{verde} y \texttt{refactor} que muestra la E8.
  \item[\texttt{docs/informe.pdf}.] Este informe, cuya fuente está en \texttt{latex/}.
\end{description}

\end{document}
